\hypertarget{appendix-f28eeec97a3d}{}
\section*{Spline bases and curve evaluation}
\section{Spline bases and curve evaluation}\label{spline-bases-and-curve-evaluation}

\textbf{Curriculum depth:} Developed · \textbf{First seen:} University courses · \textbf{Latest development:} University courses

\subsection{Definition and mechanism}\label{definition-and-mechanism}

A spline is a piecewise polynomial curve whose neighboring pieces join under explicit continuity constraints. A cubic Bézier segment uses four control points and Bernstein weights; the weights sum to one, so the curve lies in their convex hull. A B-spline uses local basis functions so moving one control point changes only nearby spans rather than the whole curve. Knots mark where polynomial pieces meet; repeated knots reduce continuity. Catmull--Rom curves interpolate their control points and can be convenient for motion paths, while Bézier curves expose endpoint tangents through their handles. To evaluate a curve, compute basis weights at parameter \(t\) and take their weighted sum of control points. Position, first derivative, and parameterization speed matter separately: even a smooth path may speed up or slow down unless its parameter is adjusted.

\subsection{Prerequisites}\label{prerequisites}

\begin{itemize}
\tightlist
\item
  \href{}{Basis coordinates and dimension (01-Foundations-I, appendix-ce0341f7ea0d)}
\end{itemize}

\subsection{Used by}\label{used-by}

\begin{itemize}
\item
  \href{Continuous\%20Notes/University\%20-\%20Worlds,\%20Images\%20and\%20Inference\%20-\%20Graphics,\%20Vision\%20and\%20Photography}{Worlds, Images and Inference}
\item
  \href{}{Bézier curves and De Casteljau evaluation (02-Foundations-II, appendix-ddbec701898c)}
\item
  \href{}{Cubic Hermite interpolation (02-Foundations-II, appendix-db3aef8f2fd6)} \#\# Full teaching explanation
\end{itemize}

\href{Continuous\%20Notes/University\%20-\%20Worlds,\%20Images\%20and\%20Inference\%20-\%20Graphics,\%20Vision\%20and\%20Photography\#Curves\%20as\%20controllable\%20functions}{Curves as controllable functions} · \href{Continuous\%20Notes/University\%20-\%20Worlds,\%20Images\%20and\%20Inference\%20-\%20Graphics,\%20Vision\%20and\%20Photography\#Animation\%20curves\%20and\%20motion\%20paths}{Animation curves and motion paths}
